\medskip

\textbf{Partie A}

\medskip

\begin{enumerate}
\item On complète le tableau des effectifs :

\begin{center}
\begin{tabularx}{0.5\linewidth}{|l|*{3}{>{\centering \arraybackslash}X|}}\cline{2-4}
\multicolumn{1}{c|}{~} & Fille & Garçon&Total \\ \hline
Anglais & 712 & 728&\np{1440} \\ \hline
Autre LV1 & 288 & 272&\np{560} \\ \hline
Total&\np{1000}&\np{1000}&\np{2000} \\ \hline
\end{tabularx}
\end{center}

Il y a donc autant de filles que de garçons.
\item 712 filles ont choisi l'Anglais, donc $p(A \cap F) = \dfrac{712}{\np{2000}}$.
\item On a $p_F(A) = \dfrac{p(F \cap A)}{p(F)} = \dfrac{p(A \cap F)}{p(F)} = \dfrac{\dfrac{712}{\np{2000}}}{\dfrac{1000}{\np{2000}}} = \dfrac{712}{\np{1000}}$.
\item On a $p(A) \times p(F) = \dfrac{712}{\np{2000}} \times \dfrac12 = \dfrac{712}{\np{4000}} = \dfrac{356}{\np{2000}}$.

Donc $p(A) \times p(F) \ne p(A \cap F)$ : les évènements $A$ et $F$ ne sont pas indépendants.
\item 728 élèves garçons ont choisi l'anglais et 272 ne l'ont pas choisi : or $728 \ne 3 \times 272$ car ce produit se termine par un 6.

L'affirmation est fausse.
\end{enumerate}

\bigskip

\textbf{Partie B}

\medskip

\begin{enumerate}
\item On a $p(P) = \dfrac{1}{4}$, donc $p(F) = 1 - p(P) = 1 - \dfrac{1}{4} = \dfrac34$.
\item~

\begin{center}
\pstree[treemode=R,nodesep=1ex,levelsep=2.5cm]{\TR{}}
{\pstree{\TR{$P~~$}\taput{\small $\frac14$}}
		{\pstree{\TR{$P~~$}\taput{\small $\frac14$}}
			{\TR{$P~~$}\taput{\small $\frac14$}
			\TR{$F~~$}\tbput{\small $\frac34$}
			}
		\pstree{\TR{$F~~$}\tbput{\small $\frac34$}}
			{\TR{$P~~$}\taput{\small $\frac14$}
			\TR{$F~~$}\tbput{\small $\frac34$}
			}
		}
\pstree{\TR{$F~~$}\tbput{\small $\frac34$}}
		{\pstree{\TR{$P~~$}\taput{\small $\frac14$}}
			{\TR{$P~~$}\taput{\small $\frac14$}
			\TR{$F~~$}\tbput{\small $\frac34$}
			}
		\pstree{\TR{$F~~$}\tbput{\small $\frac34$}}
			{\TR{$P~~$}\taput{\small $\frac14$}
			\TR{$F~~$}\tbput{\small $\frac34$}
			}
	}
}
\end{center}

\item La probabilité d'obtenir une fois Pile est :

$p(PFF) + p(FPF) + p(FFP) = \frac14 \times \frac34 \times \frac34 + \frac34 \times \frac14 \times \frac34 + \frac34 \times \frac34 \times \frac14 = \frac{9}{64} + \frac{9}{64} + \frac{9}{64} = \dfrac{27}{64}$.
\item La probabilité de ne jamais obtenir pile est égale à :

$p(FFF) = \dfrac34 \times \dfrac34 \times \dfrac34 = \dfrac{27}{64}$ (branche du bas).
\end{enumerate}
